\scanpage{180}
% L01–02
comme un quotient de $G$
\[
S\simeq G/\deleted{\unsure{G_{s_0}|A|B}}U\qquad\text{où }\deleted{\unsure{H|A|B}}U=G_{s_0}\tag*{(4)}
\]

% L03–04
ce qui identifie $Z.s_0$ à l'ensemble
\[
\operatorname{Norm}_G(\deleted{G_{s_0}}U)/\deleted{G_{s_0}}U\subset G/\deleted{G_{s_0}}U\tag*{(5)}
\]

% L05–07
$Z$ lui-même étant isomorphe canoniquement au groupe quotient
\[
Z\simeq\operatorname{Norm}_G(\deleted{G_{s_0}}U)/\deleted{G_{s_0}}U\tag*{(6)}
\]

% L08–11
l'isomorphisme étant obtenu en associant à tout $\deleted{\unsure{g}}a\in G$ normalisant $\deleted{\unsure{G_s}}U$ la ``translation à droite''
\[
\deleted{\unsure{a}}gU\longmapsto ga^{-1}U\tag*{(7)}
\]

% L12–15
(pour que $ga^{-1}U$ ne dépende que de $\inserted{\dot g=}gU\in G/U$, et non du choix de $g$, il f et s justement que \deleted{l'on ait $\operatorname{int}(a)\,U\subset$} \inserted{l'on ait} \unsure{a|int(a)}
\[
\operatorname{int}(a)\,U\subset U\tag*{(8)}
\]

% L16–21
et pour que l'application ainsi obtenue (7) de $G/U$ dans lui-$=$ (qui commute aux op. de $G$ et est surjective) soit bijective, il f. et s. que $a^{-1}$ satisfasse la $=$ condition \unsure{8)|(8)}, i.e. qu'on ait
\[
\operatorname{int}(a^{-1})\,U\deleted{\unsure{U\subset|\text{the}}} \subset U\qquad\text{qui équivaut à :}\quad U\subset\operatorname{int}(a)\,U \tag*{(9)}
\]
i.e. finalement qu'on ait compte tenu de (8) )
\[
\operatorname{int}(a)\,U=U\qquad\text{i.e. }a\in\operatorname{Norm}_G(U).\tag*{(10)}
\]

% archive sheet number: 179
